% ===== PRÉAMBULE CANONIQUE — à coller VERBATIM en tête de CHAQUE .tex =====
\documentclass[11pt,a4paper]{article}
\usepackage[utf8]{inputenc}\usepackage[T1]{fontenc}\usepackage{lmodern}
\usepackage[french]{babel}
\usepackage{amsmath,amssymb,mathtools}\usepackage{icomma}
\usepackage{siunitx}\sisetup{locale=FR}  % \qty{}{}, \si{}, \ang{}
\usepackage{xcolor}\usepackage{colortbl}
\usepackage{tikz}\usepackage{pgfplots}\pgfplotsset{compat=1.18}
\usepackage[siunitx,europeanresistors]{circuitikz} % circuits (élec.)
\usepackage[most]{tcolorbox}\usepackage{enumitem}\usepackage{array,booktabs}
\usepackage[a4paper,margin=2.2cm]{geometry}
\usetikzlibrary{arrows.meta,calc,angles,quotes,positioning,patterns,%
  backgrounds,shapes.geometric,decorations.pathmorphing,decorations.markings,intersections}
\definecolor{primary}{RGB}{222,49,99}\definecolor{primarydark}{RGB}{150,22,55}
\definecolor{defcol}{RGB}{0,114,178}\definecolor{methcol}{RGB}{52,92,148}
\definecolor{excol}{RGB}{0,135,90}\definecolor{attncol}{RGB}{200,40,55}
\tcbset{boxbase/.style={enhanced,breakable,boxrule=0.9pt,arc=2.5pt,
  left=3mm,right=3mm,top=2.3mm,bottom=2.3mm,fonttitle=\bfseries,coltitle=white}}
% --- boîtes TUTORIEL (noms EXACTS, ne pas renommer) ---
\newtcolorbox{cle}[1][]{boxbase,colback=defcol!6,colframe=defcol,
  title={\textsc{À retenir}\ifx\relax#1\relax\relax\else\textmd{ : \itshape #1}\fi}}
\newtcolorbox{methode}[1][]{boxbase,colback=methcol!7,colframe=methcol,sharp corners,
  title={\textsc{Méthode}\ifx\relax#1\relax\relax\else\textmd{ --- \itshape #1}\fi}}
\newtcolorbox{exemplebox}[1][]{boxbase,colback=excol!5,colframe=excol,
  title={\textsc{Exemple}\ifx\relax#1\relax\relax\else\textmd{ : \itshape #1}\fi}}
\newtcolorbox{piege}[1][]{boxbase,colback=attncol!6,colframe=attncol,
  title={\textsc{Piège}\ifx\relax#1\relax\relax\else\textmd{ : \itshape #1}\fi}}
\newtcolorbox{remarque}[1][]{boxbase,colback=black!4,colframe=black!45,
  title={\textsc{Remarque}\ifx\relax#1\relax\relax\else\textmd{ : \itshape #1}\fi}}
% --- boîtes EXERCICE / CORRIGÉ (noms EXACTS, auto-comptées) ---
\newtcolorbox[auto counter]{exo}[1][]{boxbase,colback=primary!4,colframe=primary,
  title={\textsc{Exercice}~\thetcbcounter\ifx\relax#1\relax\relax\else\textmd{ --- \itshape #1}\fi}}
\newtcolorbox[auto counter]{corrige}[1][]{boxbase,colback=excol!4,colframe=excol,
  title={\textsc{Corrigé}~\thetcbcounter\ifx\relax#1\relax\relax\else\textmd{ --- \itshape #1}\fi}}
\newcommand{\vv}[1]{\vec{#1}}\newcommand{\ds}{\displaystyle}
\newcommand{\R}{\mathbb{R}}\newcommand{\N}{\mathbb{N}}\newcommand{\Z}{\mathbb{Z}}
% (un \usepackage{fancyhdr}+en-tête est possible ; il est ignoré côté web)
% =========================================================================
\begin{document}
% THEME_TITLE: Mouvements rectilignes : MRU, MRUA/MRUD et Torricelli
\sisetup{per-mode=symbol}

\noindent Sur une droite, tout mouvement se ramène à deux familles : à \textbf{vitesse constante}
(MRU) ou à \textbf{accélération constante} (MRUA/MRUD). Trois équations et une relation « sans le temps » suffisent.

\begin{cle}[mouvement rectiligne uniforme (MRU)]
Vitesse constante, accélération nulle :
\[ a=0,\qquad v=v_0=\text{cte},\qquad x(t)=x_0+v_0\,t. \]
Graphiquement : $v(t)$ est une horizontale, $x(t)$ une droite de pente $v_0$.
\end{cle}

\begin{cle}[mouvement uniformément accéléré / décéléré (MRUA/MRUD)]
Accélération constante :
\[ a=\text{cte},\qquad v(t)=v_0+a\,t,\qquad x(t)=x_0+v_0\,t+\tfrac12 a\,t^2. \]
$a>0$ : mouvement accéléré (MRUA) ; $a<0$ : décéléré (MRUD, freinage).
Graphiquement : $v(t)$ est une droite de pente $a$, $x(t)$ une parabole.
\end{cle}

\begin{cle}[relation de Torricelli (sans le temps)]
En éliminant $t$ entre $v=v_0+at$ et l'équation horaire :
\[ \boxed{\,v^2=v_0^2+2\,a\,(x-x_0)\,}. \]
Idéale quand la \emph{durée est inconnue} (distance de freinage, par exemple).
\end{cle}

\begin{methode}[quelle équation choisir ?]
\begin{center}\small
\begin{tabular}{@{}>{\raggedright\arraybackslash}p{5.4cm}|>{\raggedright\arraybackslash}p{8.6cm}@{}}
\toprule
\rowcolor{primary!12}\textbf{Ce qu'on cherche / connaît} & \textbf{Équation à utiliser} \\
\midrule
On dispose du temps $t$ & $v=v_0+at$ \quad ou \quad $x=x_0+v_0t+\tfrac12at^2$ \\
Le temps n'apparaît pas (vitesses \& distances) & Torricelli : $v^2=v_0^2+2a\,\Delta x$ \\
Vitesse constante & MRU : $x=x_0+v_0t$ \\
\bottomrule
\end{tabular}
\end{center}
\end{methode}

\begin{exemplebox}[distance de freinage]
Une voiture à $v_0=\qty{20}{\metre\per\second}$ freine à $a=\qty{-5}{\metre\per\second\squared}$.
Torricelli ($v=0$) : $\Delta x=\dfrac{-v_0^2}{2a}=\dfrac{-400}{-10}=\qty{40}{\metre}$ ; durée $t=\dfrac{-v_0}{a}=\qty{4}{\second}$.
\end{exemplebox}

\begin{piege}[la distance de freinage varie comme le carré de la vitesse]
Comme $\Delta x=\dfrac{v_0^2}{2|a|}$, \textbf{doubler} la vitesse initiale \textbf{quadruple} la distance
d'arrêt (à freinage égal). Passer de \qty{50}{} à \qty{100}{\kilo\metre\per\hour} multiplie par 4 la distance !
Pense aussi à garder le \emph{signe de $a$} cohérent avec l'orientation de l'axe.
\end{piege}

\begin{remarque}[tableau récapitulatif]
\begin{center}\small
\begin{tabular}{@{}l|c|c|c@{}}
\toprule
\rowcolor{primary!12}\textbf{Mouvement} & \textbf{Accélération} & \textbf{Vitesse} & \textbf{Position} \\
\midrule
MRU & $a=0$ & $v=v_0$ & $x=x_0+v_0t$ \\
MRUA / MRUD & $a=\text{cte}$ & $v=v_0+at$ & $x=x_0+v_0t+\tfrac12at^2$ \\
\bottomrule
\end{tabular}
\end{center}
\end{remarque}

\end{document}
